\chapter{確率・情報・エントロピー}

\lectureoverview{統計力学は、微視的法則と不完全な知識とが出会うところから始まる。まず確率を精密に定式化し、次に決定論的運動それ自体がどのように確率を割り当てるのか、保存則がどのように状態空間を分割するのか、そして情報を保存する力学がなぜ自然にエントロピーへと導くのかを問う。第一法則と第零法則が入ってくるのは、この構造が整ってからである。}

\section{なぜ統計力学か}

統計力学はときに古い物理学だと言われる。その説明は的を外している。統計力学は前近代的であり、現代的であり、そして我々が知る理論に取って代わるいかなるポストモダン物理学においても有用であり続けるだろう。個別のモデルは現れては消えるが、熱力学第二法則は繰り返し道標として生き残ってきた。自然についての提案された記述が、秩序をただで作り出したり、情報を根本的に消去したり、閉じた系を不自然に単純な状態へと走らせたりするように見えるときはいつでも、第二法則は何かが取りこぼされていると警告する。

この主題は、新たに発見された粒子ほど華やかには見えないかもしれないが、少なくとも同じくらい深く、はるかに一般的である。気体、液体、固体、放射、ブラックホール、情報、量子場は、いずれも統計的推論を必要とする。実際、粒子を発見するために用いられた実験装置も、それなしには理解できない。何世代もの優れた物理学者が、主たる仕事が別のところにあっても統計力学を習得したのはこのためである。この主題は有用であるだけでなく、際立って美しい。

なぜそれが必要なのかを見るには、力学が与える理想から始めればよい。閉じた系とは、考察の対象となっている宇宙全体か、あるいは周囲が影響を及ぼさないほど十分に孤立した部分系である。その初期条件とその発展法則を知っていれば、古典力学は完全な予測を与える。量子力学は予測できるものを変えるが、それもまた精密に定義された状態と、理論が許す最大限の情報を予測する法則とを持つ。

原理上の予測と有用な理解との間には厳しい違いがある。この部屋のすべての分子の位置と速度の一覧は、途方もなく長く、ほとんど即座に変化し、我々が実際に気にかける問いのほとんどに答えてくれない。したがって、正確な微視的法則は、完全に予測的でありながら同時に実用上は役に立たないことがありうる。

統計力学は、この物理的状況に適用された確率論である。初期条件が完全な精度で知られていないとき、法則が近似的にしか知られていないとき、系が開いているとき、あるいは単に微視的変数の数が多すぎて追跡できないときに、それは必要となる。これらのいずれも予測が弱いことを意味しない。大きな数はしばしば統計的予測を極めて鋭いものにする。

孤立した気体の箱を考えよう。我々はすべての分子の軌道を予測することはできないが、わずかな巨視的データから、圧力、エネルギー、その他多くのバルク量を高い精度で予測できる。不確かなまま残るのは揺らぎである。予測できない瞬間に、ある小さな領域が平均より多くの分子を含み、別の領域がより少なく含むことがある。そのような事象は確率分布の裾に位置する。起こりにくいが、禁じられてはいない。特定の発生の瞬間を予測することなく、そのような揺らぎがどのくらいの頻度で起こるべきかを計算できる。

硬貨投げは同じ区別を示す。十億回の公正な投げでは、表と裏の数はほぼ等しくなる。十分に多く投げれば、表が長く連続することもいつかは起こるが、確率論が予測するのはその頻度であって、あらかじめその時刻を指定することではない。統計力学は、その要素が小さすぎるか、多すぎるか、あるいはその両方である系のために設計されている。

本講義では、法則そのものと物質への応用の両方を展開する。この主題は繰り返し学ぶことに報いてくれる。一度学び、その一部を忘れ、どの考えが基本的であるかについてより深い感覚をもって再び学ぶのである。我々は、続くすべてを支配するその考えから始める。

\section{原始的な言語としての確率}

なぜ確率がうまく働くのかは、確率の計算それ自体が説明するものではない。確率的予測は特定の試行では外れることがある。結果が可能ではあるが起こりにくいとは、そういうことである。理論は例外がどのように分布するかを教える。例外を廃止するのではない。したがって我々は確率を原始的概念として受け入れ、それが従う規則を述べる。

\subsection{状態と規格化}

\(i\) が可能な結果、または系の可能な状態を表すとしよう。1個の硬貨では、\(i\) は表または裏を意味しうる。1個のサイコロでは、\(1\) から \(6\) までの値をとりうる。2個のサイコロでは、添字の対が両方の面を記録する。まず有限の離散的な状態空間から始める。
\begin{equation}
 i=1,\ldots,N,
\end{equation}
そして後に、無限個および連続無限個の状態へ移行する。

無知であるとは、系が状態をもたないという意味ではない。どの状態をもつかを我々が知らないという意味であり、それゆえ各可能性に確率 \(p_i\) を割り当てる。基本規則は
\begin{align}
 p_i &\geq 0, \label{eq:sm1-prob-positive}\\
 \sum_{i=1}^{N}p_i &=1. \label{eq:sm1-prob-normalization}
\end{align}
である。諸選択肢は互いに排反かつ網羅的である。すなわち、そのうちのちょうど1つが生じる。負の確率はここでは役割をもたず、和が1であるのは、許された結果のどれかが必ず生じるからである。

\subsection{頻度と大数の法則}

確率は繰り返し試行によって実験と結びつく。同じ実験を \(N_{\mathrm{trials}}\) 回行い、\(N(i)\) を結果 \(i\) を生じる試行の数としよう。このとき頻度解釈は
\begin{equation}
 p_i=
 \lim_{N_{\mathrm{trials}}\to\infty}
 \frac{N(i)}{N_{\mathrm{trials}}}.
 \label{eq:sm1-frequency}
\end{equation}
である。これは理想的な極限を通じて表された物理的仮説である。無限回の実験を行う実験室は存在しない。主張は、十分に多くの試行を行えば、その比が意図された予測にとって十分近く \(p_i\) に落ち着くというものである。有限の標本は極限値のまわりで変動する。

\subsection{平均値}

次に、各状態に数 \(f(i)\) を割り当てる。それはエネルギー、運動量、あるいは我々が測定したい任意の他の量でありうる。その期待値は
\begin{equation}
 \langle f\rangle=\sum_i f(i)p_i.
 \label{eq:sm1-expectation}
\end{equation}
である。ブラケット記法は量子力学の外でも便利である。頻度の言葉では、同じ平均は
\begin{equation}
 \langle f\rangle
 =
 \lim_{N_{\mathrm{trials}}\to\infty}
 \frac{1}{N_{\mathrm{trials}}}
 \sum_i f(i)N(i).
 \label{eq:sm1-frequency-average}
\end{equation}
である。

硬貨について、
\begin{equation}
 f(\mathrm{H})=+1,
 \qquad
 f(\mathrm{T})=-1.
\end{equation}
と定義する。表と裏がそれぞれ確率 \(1/2\) をもつならば、
\begin{equation}
 \langle f\rangle
 =(+1)\frac12+(-1)\frac12=0.
\end{equation}
である。平均は可能な結果である必要はない。硬貨が0の上に落ちることは決してない。0はその2つの可能な値の分布を要約しているのである。

この短い間奏は、この分野全体で用いられる語彙を確定する。すなわち、状態、確率、頻度、そして期待値である。

\section{確率を割り当てる3つの方法}

\subsection{対称性}

なぜ我々は通常、硬貨の2つの面に等しい確率を割り当てるのか。第一の答えは対称性である。ほぼ理想的な硬貨を裏返すことは、実験を変えることなく表と裏を交換する。そのとき、どちらかの結果を優先する物理的理由はない。実際の刻印や不完全さはわずかな偏りをもたらすが、対称な模型は意図された実験を捉えている。
\begin{equation}
 p(\mathrm{H})=p(\mathrm{T})=\frac12.
\end{equation}

同じ議論により、対称なサイコロの各面に確率 \(1/6\) が与えられる。より一般には、対称性によって関係づけられる配置には、何らかの追加情報がその対称性を破らない限り、等しい確率が割り当てられる。

\subsection{実験}

歪んだサイコロには、6つの面すべてを交換する対称性はない。その場合、対称性だけでは答えは得られない。偏りを予測するより深い力学的模型をもつことはありうるが、そのような模型がなければ、誠実な手続きは実験的である。すなわち、サイコロを多数回投げ、頻度を記録し、それらが安定するのを待ち、得られた確率の表を後の予測に用いるのである。

対称性と実験は競合する定義ではない。それらは確率分布に関する情報の2つの源である。統計力学は対称性に大きく依拠することになるが、力学が第三の道を提供する。

\subsection{決定論的運動からの確率}

6つの色の配置をもつ不規則な物体を想像しよう。それは、確定した更新則に従って、固定された離散的な時間間隔で状態を変える。
\begin{equation}
 R\longrightarrow B\longrightarrow Y\longrightarrow G
 \longrightarrow O\longrightarrow P\longrightarrow R.
 \label{eq:sm1-six-cycle}
\end{equation}
物体自体は幾何学的対称性をもつ必要はない。この法則は、与えられた配置にどの配置が続くかを述べているだけである。

\begin{figure}[htbp]
\centering
\includegraphics[width=0.92\textwidth]{lecture_01_figure_02.png}
\caption{期待値の公式は左の板に残ったままで、6つの決定論的な色の更新が右に列挙されている。}
\lectureframe{00:26:28}
\end{figure}

板のリストは、有向サイクルとして描くとより明瞭である。
\begin{figure}[htbp]
\centering
\begin{tikzpicture}[
  >=Latex,
  state/.style={draw,circle,minimum size=8mm,inner sep=1pt}
]
\node[state] (R) at (90:2.2) {$R$};
\node[state] (B) at (30:2.2) {$B$};
\node[state] (Y) at (-30:2.2) {$Y$};
\node[state] (G) at (-90:2.2) {$G$};
\node[state] (O) at (-150:2.2) {$O$};
\node[state] (P) at (150:2.2) {$P$};
\draw[->] (R) -- (B);
\draw[->] (B) -- (Y);
\draw[->] (Y) -- (G);
\draw[->] (G) -- (O);
\draw[->] (O) -- (P);
\draw[->] (P) -- (R);
\end{tikzpicture}
\caption{単一の6状態軌道。各周期の間にすべての状態が1回ずつ訪れられる。}
\label{fig:sm1-six-cycle}
\end{figure}

初期状態が分からないものとし、未知の時刻にスナップショットを撮るとする。各段階が等しい時間だけ続くので、系は各色に一周期の六分の一ずつを費やす：
\begin{equation}
 p(R)=p(B)=p(Y)=p(G)=p(O)=p(P)=\frac16.
 \label{eq:sm1-uniform-cycle}
\end{equation}
この等式は物体の形に由来するものではない。各状態に費やす時間が等しいことに由来する。

矢印を並べ替えて、依然として1つの6状態サイクルを形成しながら、異なる微視的法則を得ることができる。詳細な順序は変わるが、各状態が依然として周期ごとに1回ずつ生じるので、スナップショットの確率は \(1/6\) のままである。したがってこの予測は、出発点も詳細な順序も必要としない。必要とするのは、軌道構造と更新あたりの等しい時間だけである。

\begin{classroomqa}
\classquestion{対象が対称でない場合、すべての色に確率\(1/6\)を与えることはどのような意味で正当なのか。}
\lecturerresponse{その割り当ては幾何学的というより動力学的である。6段階の軌道の未知の位相で撮られたスナップショットは、6つの等しい時間スロットをサンプリングする。関係する等価性は、対象の面の間ではなく、軌道上の位置の間にある。滞留時間が等しくなければ、確率も等しくならない。
\lecturetimestamp{00:27:11--00:29:49}}
\end{classroomqa}

\section{非連結なサイクルと保存量}

決定論的法則は、すべての状態を1つの軌道に置く必要はない。2つの非連結なサイクルを考える。
\begin{equation}
 R\longrightarrow B\longrightarrow G\longrightarrow R,
 \qquad
 P\longrightarrow Y\longrightarrow O\longrightarrow P.
 \label{eq:sm1-two-cycles}
\end{equation}
系が第1のサイクル上にあるとわかっているなら、
\begin{equation}
 \begin{aligned}
 p(R)=p(B)=p(G)&=\frac13,\\
 p(P)=p(Y)=p(O)&=0.
 \end{aligned}
\end{equation}
第2のサイクル上にあるという知識は、これらの割り当てを逆にする。

より一般的な場合は事前重みを必要とする。サイクルを\(+1\)と\(-1\)とラベル付けし、それらの確率を\(p_+\)と\(p_-\)とする。ただし
\begin{equation}
 p_++p_-=1.
\end{equation}
このとき
\begin{align}
 p(R)=p(B)=p(G)&=\frac{p_+}{3},\\
 p(P)=p(Y)=p(O)&=\frac{p_-}{3}.
 \label{eq:sm1-sector-weights}
\end{align}
因子\(1/3\)はサイクル内の運動から従う。数\(p_+\)と\(p_-\)はそうではない。それらは追加の知識、おそらく初期セクターを選択する別の実験または別の系から来なければならない。

ラベル\(+1\)と\(-1\)は保存量を定義する。いったん系が一方のサイクルで始まれば、更新則がそれを他方へ運ぶことは決してない。この量をエネルギーと呼んでもよいし、\emph{zilch}のような仮の名前を与えてもよい。その名前は二次的である。保存とは、状態空間が力学の越えられない領域に分割されることを意味する。

サイクルは同じ数の状態を含む必要はない。4状態の軌道と2状態の軌道も依然として2つの保存セクターを定義する。このとき等時間サンプリングは、第1のセクター内で\(1/4\)を、第2のセクター内で\(1/2\)を与え、それぞれの場合にそのセクターの事前確率を掛ける。統計力学は論理的に別個の2つの問いに答えなければならない。
\begin{enumerate}
\item 与えられた保存セクターにおいて、力学はその状態にどれだけの時間を費やすか。
\item 保存セクター自体にはどのような確率を割り当てるべきか。
\end{enumerate}

\begin{classroomqa}
\audiencequestion{1つの保存セクター内のすべての状態が同じ確率を受け取ると仮定しているのか。}
\lecturerresponse{ここに示した離散的な玩具法則については、はい。各状態が同じ時間間隔続き、運動が決定論的だからである。無知は軌道の初期位相にある。ミリ秒の時間分解能をもつ検出器はマイクロ秒の更新を分解できないが、十分に速いスナップショットは捉えられた状態を依然として同定できる。より一般的な系では、等しくない滞留時間を明示的に含めなければならない。
\lecturetimestamp{00:36:16--00:37:17}}
\end{classroomqa}

\begin{classroomqa}
\audiencequestion{2つのスナップショットが撮られ、最初のものが一方のサイクルを同定したなら、2番目のものも同じサイクルを同定すべきか。}
\lecturerresponse{はい。まさにそれが保存の意味である。いったん測定が保存量の値を決定すれば、将来の確率はその値に条件付けることができ、他方のセクターの状態は確率ゼロを受け取る。
\lecturetimestamp{00:37:18--00:37:45}}
\end{classroomqa}

\subsection{2つの系の間で共有されるエネルギー}

ここで抽象的なラベルをエネルギーに置き換える。1つの孤立した箱は閉じた系であり、保存されたエネルギーをもつ。2つの孤立した箱は、相互作用しないため、2つの別々に保存されたエネルギーをもつ。エネルギーが流れることを許す狭い通路でそれらを接続すると、全エネルギーは固定されたまま、別々のエネルギーは変化しうる。

\begin{figure}[htbp]
\centering
\includegraphics[width=0.92\textwidth]{lecture_01_figure_03.png}
\caption{エネルギーが通過できる狭いチャネルで接続された2つの部分系。}
\lectureframe{00:39:19}
\end{figure}

\begin{figure}[htbp]
\centering
\begin{tikzpicture}[>=Latex]
\draw[thick,rounded corners=4pt] (-4,1.8) rectangle (-0.7,-1.5);
\draw[thick,rounded corners=4pt] (0.7,1.5) rectangle (4,-1.8);
\draw[thick] (-0.7,0.45) -- (0.7,0.45);
\draw[thick] (-0.7,-0.35) -- (0.7,-0.35);
\draw[<->,very thick] (-0.45,0.05) -- (0.45,0.05);
\node at (-2.35,0.15) {subsystem \(1\)};
\node at (2.35,-0.15) {subsystem \(2\)};
\node[above] at (0,0.5) {energy exchange};
\end{tikzpicture}
\caption{接続が開かれると、\(E_1+E_2\)だけが保存される。エネルギーの分割は確率の問題になる。}
\label{fig:sm1-connected-systems}
\end{figure}

全エネルギーが与えられたとき、部分系\(1\)が\(E_1\)をもち、部分系\(2\)が残りをもつ確率を問うことができる。等しい箱は平均して等しいエネルギーをもつべきだが、平均は分布の全体ではない。したがって部分系の確率を決定する1つの方法は、それをより大きな閉じた系に埋め込み、その保存則が欠けている情報を与えるようにすることである。この構成は統計力学の至るところで再び現れる。

\section{良い法則、悪い法則、そしてマイナス第一法則}

すべての決定論的更新則が基礎法則として許されるわけではない。写像
\begin{equation}
 \begin{aligned}
 R&\mapsto R,& B&\mapsto R,& Y&\mapsto R,\\
 G&\mapsto R,& O&\mapsto R,& P&\mapsto R.
 \end{aligned}
 \label{eq:sm1-all-to-red}
\end{equation}
を考える。これは未来を完全に予測する。1ステップ後、状態は赤である。しかし逆方向では失敗する。系が赤になると、青、黄、緑、橙、紫、または赤自身から来たのかを何も明かさない。

\begin{figure}[htbp]
\centering
\begin{tikzpicture}[
  >=Latex,
  state/.style={draw,circle,minimum size=8mm,inner sep=1pt}
]
\node[state,fill=red!12] (R) at (0,0) {$R$};
\node[state] (B) at (-3,1.6) {$B$};
\node[state] (Y) at (-3,0) {$Y$};
\node[state] (G) at (-3,-1.6) {$G$};
\node[state] (O) at (3,1.6) {$O$};
\node[state] (P) at (3,-1.6) {$P$};
\draw[->] (B) -- (R);
\draw[->] (Y) -- (R);
\draw[->] (G) -- (R);
\draw[->] (O) -- (R);
\draw[->] (P) -- (R);
\draw[->,loop above] (R) to (R);
\end{tikzpicture}
\caption{多対一の法則は次の状態を予測するが、可能な過去の状態の間の区別を消去する。}
\label{fig:sm1-bad-law}
\end{figure}

対照的に、閉じたサイクルは記録を保存する。経過時間がわかっていれば、何千周期も経た後でさえ、状態を未来へも過去へも発展させることができる。法則は初期条件の間の区別を保存する。Susskindはこの要請を\emph{物理学の第一法則に対するマイナス}と呼ぶ。情報は根本的には失われないのである。

有限の離散状態空間に対しては、図式的な判定基準は単純である。
\begin{quote}
良い法則の各状態は、ちょうど1本の出ていく矢印と、ちょうど1本の入ってくる矢印をもつ。
\end{quote}
出ていく矢印が未来を決定し、入ってくる矢印が過去を再構成する。数学的には、更新は状態の置換である。

\subsection{摩擦は見かけ上の反例にすぎない}

連続的な対応物はLiouvilleの定理である。摩擦はこれに矛盾するように見える。消しゴムを机の上で滑らせると、消しゴムがどこから出発しても、やがて静止する。模式的な抗力方程式は
\begin{equation}
 m\ddot x=-\gamma\dot x,
 \qquad \gamma>0.
 \label{eq:sm1-drag}
\end{equation}
であり、その速度は指数関数的に減衰する。もしこれが密閉された部屋の中のあらゆる粒子の根本法則であるなら、すべての運動は消え去り、温度はゼロまで下がり、多くの初期状態が1つの最終状態へと崩壊することになる。エネルギー保存は破れ、情報は消え、無秩序は減少する。これは、すべてが赤へ向かうことの連続版であろう。

実際の摩擦はこのようにはならない。消しゴムは机と、それ自身と、周囲の環境とを温める。微視的な位置と運動量が、可視的な運動から失われた区別を記録する。摩擦は、観測されない自由度が省略された後の部分系に対する有効方程式であり、完全な閉じた系に対する根本的な散逸法則ではない。

\begin{classroomqa}
\classquestion{微分方程式
\(m\ddot x=-\gamma\dot x\)は何を表し、なぜ根本的な閉じた系の法則ではないのか。}
\lecturerresponse{それは粘性抗力の方程式である。それ自体ではあらゆる速度をゼロへと駆り立て、初期温度にかかわらず気体をゼロまで冷却させることになる。実際の閉じた記述では、失われた力学的エネルギーと情報は微視的な自由度へと移る。抗力方程式は、それらの変数が抑圧された後にのみ信頼できる。
\lecturetimestamp{00:44:40--00:48:56}}
\end{classroomqa}

\section{占有状態と最初のエントロピー}

情報保存には等価な確率的な言明がある。\(N\)個の状態のうち\(M\)個に等しい確率が割り当てられているとしよう。
\begin{equation}
 p_i=
 \begin{cases}
 1/M, & i\text{ in the occupied set},\\
 0, & i\text{ outside it}.
 \end{cases}
 \label{eq:sm1-uniform-occupied}
\end{equation}
情報を保存する更新は、どの状態が占有されているかを入れ替えることはあっても、それらを融合させることはできない。それ以降のどの時刻でも、依然として確率\(1/M\)をもつ\(M\)個の状態がある。これに対して、悪い写像\eqref{eq:sm1-all-to-red}は、多くの占有状態を1つに崩壊させうる。

数\(M\)は無知を測る。\(M=N\)ならば、あらゆる状態が可能であり、我々の無知は最大である。\(M=1\)ならば、微視的状態は正確にわかっている。したがってエントロピーの最初の定義は
\begin{equation}
 S=\log M.
 \label{eq:sm1-entropy-logm}
\end{equation}
であり、これは
\begin{equation}
 S_{\max}=\log N,
 \qquad
 S_{\min}=0.
\end{equation}
を与える。状態空間が無限であれば、有限の上限が存在する必要はない。

この展開では、エントロピーは温度より前に現れ、エネルギーの保存の形式的な言明よりもさらに前に現れる。その意味でそれはより原始的である。しかしそれは力学的系だけの性質ではない。それは系と、何が知られているかを表す確率分布とに共同で依存する。同じ力学的系でも、異なる情報をもつ観測者によって異なるエントロピーを割り当てられうる。

厳密な微視的発展の下では、占有状態は並べ替えられるだけなので、この細粒エントロピーは保存される。実際には我々はタイミングを失い、相関を捨て、膨大な数の変数を追跡し損なう。すると我々の確率分布はより多くの選択肢に広がり、粗い記述に割り当てられたエントロピーは増大する。我々の使える記述が情報を失うために、方程式が情報を破壊する必要はない。

\begin{classroomqa}
\classquestion{\(N\)個の状態をもつ有限系の最大可能エントロピーは何か。}
\lecturerresponse{それは\(\log N\)であり、\(N\)個の状態がすべて等確率のときに達成される。最小のエントロピーはゼロであり、1つの状態が確率1をもつときに達成される。したがってエントロピーは、利用可能な状態空間だけでなく、知識の状態をも測る。
\lecturetimestamp{00:56:46--00:58:19}}
\end{classroomqa}

\section{連続位相空間におけるLiouvilleの流れ}

連続的な位置と運動量をもつ粒子の場合、微視的状態は相空間の一点である。座標を\(q_a\)、その共役運動量を\(p_a\)と書く。粒子が多数になると軸の数は膨大になり、一つの\(q\)と一つの\(p\)をもつ二次元の図はあくまで代用にすぎない。すべての座標には運動量の相棒が存在するため、相空間は偶数次元である。

我々の知識が領域\(\Omega_0\)によって表されるとしよう。確率密度はその領域の内部で一様であり、外部ではゼロである。例えば、すべての粒子が部屋の中にあるという知識は位置座標を制限し、それらの速さに対する上限は運動量座標を制限する。

ハミルトン的発展は各相空間点を軌道に沿って運ぶ。近傍の点は近傍の点へと移り、占有された領域は非圧縮性流体の一片のように流れる。
\begin{figure}[htbp]
\centering
\begin{tikzpicture}[>=Latex,scale=0.95]
\draw[->] (-4.2,0) -- (4.3,0) node[right] {$q$};
\draw[->] (-3.7,-2.2) -- (-3.7,2.4) node[above] {$p$};
\draw[fill=blue!12,thick]
  (-3.0,-0.5) .. controls (-2.7,1.0) and (-1.3,1.2) .. (-0.9,0.1)
  .. controls (-0.6,-1.0) and (-2.2,-1.4) .. cycle;
\node at (-1.9,0.0) {$\Omega_0$};
\draw[->,very thick] (-0.3,0.2) -- (0.8,0.2);
\draw[fill=orange!15,thick]
  (1.2,-1.1) .. controls (1.5,0.3) and (2.3,1.8) .. (3.2,1.1)
  .. controls (3.6,0.7) and (2.8,-0.4) .. (2.5,-1.5)
  .. controls (2.1,-1.8) and (1.5,-1.6) .. cycle;
\node at (2.35,0.1) {$\Omega_t$};
\end{tikzpicture}
\caption{ハミルトン的流れは占有された相空間領域を伸長し剪断しうるが、その体積は保存する。}
\label{fig:sm1-liouville-flow}
\end{figure}

リウヴィルの定理は
\begin{equation}
 \operatorname{Vol}(\Omega_t)
 =
 \operatorname{Vol}(\Omega_0).
 \label{eq:sm1-liouville-volume}
\end{equation}
と述べる。初期密度が\(\Omega_0\)上で一様ならば、輸送された密度は\(\Omega_t\)上で一様である。これは占有された離散状態の数と重みを保存することの連続的類似物である。

この定理はまた、摩擦のパズルを幾何学的に解決する。消しゴムの可視的な運動量のみに注目すると、可能性の雲は\(p=0\)へと絞られていくように見える。完全な相空間では、その雲は消しゴムと机の分子の位置と運動量へと広がる。射影は収縮しうるが、完全な高次元体積は一定のままである。

\begin{classroomqa}
\classquestion{すべての可視的な軌道が静止に近づくならば、占有された相空間領域は軸\(p=0\)上に崩壊したのか。}
\lecturerresponse{縮約された射影においてのみである。完全な系は机、消しゴム、そして周囲の微視的変数を含む。可視的な運動量の範囲が狭まるにつれて、分布はそれらの隠れた方向へと広がる。リウヴィルの定理は完全な閉じた相空間に適用され、巨視的変数の不完全な集合には適用されない。
\lecturetimestamp{01:04:37--01:06:29}}
\end{classroomqa}

\section{マイナス第一法則から第一法則へ}

\subsection{第零法則の簡単な考察}

第一法則を述べる前に、第零法則に名前を与えておくと有用である。熱平衡はまだ定義されていないので、今のところこの法則は構造的な言明にすぎない。
\begin{quote}
系\(A\)が\(B\)と熱平衡にあり、\(B\)が\(C\)と熱平衡にあるならば、\(A\)は\(C\)と熱平衡にある。
\end{quote}
この推移性が後に温度を可能にする。今はこれを脇に置く。

\subsection{エネルギー保存}

閉じた系に対する第一法則は単に
\begin{equation}
 \frac{dE}{dt}=0.
 \label{eq:sm1-first-law}
\end{equation}
である。その簡潔な形はその力を隠している。エネルギーは保存量であり、エネルギーという語は後にさらに多くの熱力学的内容を獲得するであろう。

閉じた系が二つの部分からなり、それらのエネルギーが加法的であるならば、
\begin{equation}
 E=E_1+E_2,
\end{equation}
したがって
\begin{equation}
 \frac{dE_1}{dt}+\frac{dE_2}{dt}=0,
 \qquad\text{or}\qquad
 \frac{dE_1}{dt}=-\frac{dE_2}{dt}.
 \label{eq:sm1-energy-exchange}
\end{equation}
第二の形は交換を明示する。一方の部分系が失ったエネルギーは他方が得る。

加法性の仮定には注意が必要である。二つの重力を及ぼし合う粒子は個々の運動エネルギーを持つが、それらの相互作用はその対に属するポテンシャルエネルギーに寄与する。全エネルギーは保存されるが、それを一方の物体のエネルギーと他方のエネルギーの和として厳密に書くことはできない。

熱力学ではしばしば優れた加法的近似が許される。なぜなら、バルクのエネルギーは体積に比例してスケールする一方、大きな隣接部分間の短距離相互作用エネルギーは主にそれらの接触面積に比例してスケールするからである。大きな机をブロックに分割しよう。各ブロックの内部エネルギーは体積効果であり、隣接ブロック間の結合は表面効果である。大きなブロックでは、後者は比較的小さい。式
\eqref{eq:sm1-first-law}は閉じた全体に対して厳密であり、部分系の形
\eqref{eq:sm1-energy-exchange}は相互作用エネルギーが無視できることを仮定している。

\begin{classroomqa}
\audiencequestion{規格化\(\sum_i p_i=1\)も、列挙された結果が互いに排他的であることを仮定しているのか。}
\lecturerresponse{はい。状態ラベルは選択肢が重ならないように定義される。すなわち、サイコロが黄色の状態にあるならば、同時に赤の状態にはない。通常の言葉では色は混ざりうるが、確率計算で用いられる状態分割は互いに排他的な選択肢からなっていなければならない。
\lecturetimestamp{01:12:14--01:12:55}}
\end{classroomqa}

\section{エントロピーへ戻る}

講義はここで未完了の課題へ明示的に戻る。温度は感覚には馴染み深いが、導出された量である。温度を清潔に定義する前に、エントロピーとエネルギーが整っていなければならない。

\subsection{一般的な確率分布}

特別な定義\(S=\log M\)は、\(M\)個の状態が等確率であり、他のすべてが確率ゼロである場合に適用される。一般的な分布は矩形である必要はない。その重みは状態ラベルにわたって滑らかに、あるいは不規則に変化しうる。

\begin{figure}[htbp]
\centering
\includegraphics[width=0.92\textwidth]{lecture_01_figure_05.png}
\caption{離散的な状態ラベル上に描かれた確率分布。エントロピーに必要な特徴は、正確な曲線ではなく、広がりである。}
\lectureframe{01:14:40}
\end{figure}

\begin{figure}[htbp]
\centering
\resizebox{\linewidth}{!}{%
\begin{tikzpicture}[>=Latex,scale=0.93]
\draw[->] (-4.2,0) -- (4.4,0) node[right] {state label \(i\)};
\draw[->] (-4,0) -- (-4,3.2) node[above] {$p_i$};
\foreach \x in {-3.6,-3.1,...,3.6}
  \draw (\x,0.08) -- (\x,-0.08);
\draw[very thick]
  plot[smooth] coordinates {
  (-3.7,0.12) (-3.0,0.28) (-2.3,0.72) (-1.6,1.45)
  (-0.9,2.25) (-0.2,2.65) (0.5,2.28) (1.2,1.55)
  (1.9,0.88) (2.7,0.38) (3.5,0.14)};
\draw (-1.8,-0.18) -- (-1.8,-0.5);
\draw (1.8,-0.18) -- (1.8,-0.5);
\draw[<->] (-1.8,-0.46) -- (1.8,-0.46);
\node at (0,-0.82) {effective spread};
\end{tikzpicture}
}
\caption{狭い分布は広い分布よりも不確かさが小さい。エントロピーは分布を全体として測らなければならない。}
\label{fig:sm1-probability-spread}
\end{figure}

一般的な定義は
\begin{equation}
 S=-\sum_i p_i\log p_i.
 \label{eq:sm1-shannon-entropy}
\end{equation}
である。\eqref{eq:sm1-expectation}を用いると、これは
\begin{equation}
 S=-\langle\log p\rangle.
\end{equation}
となる。負号は、\(0<p_i<1\)のとき\(\log p_i<0\)であるという事実を補償する。

\subsection{一様な結果の回復}

一般式が最初の定義を再現することを確認しなければならない。空の状態では、\(p_i=0\)である一方、\(\log p_i\)は発散する。その積はよく定義された極限値
\begin{equation}
 \lim_{p\to0^+}p\log p=0.
 \label{eq:sm1-plogp-limit}
\end{equation}
を持つ。したがって、確率ゼロの状態は何も寄与しない。

\begin{classroomqa}
\audiencequestion{もし\(p_i=0\)ならば、その項は本当に消えるのか？ゼロの対数は負の無限大である。}
\lecturerresponse{その項は極限
\(\lim_{p\to0^+}p\log p\)によって定義され、これはゼロである。因子\(p\)は対数が発散するよりも速くゼロに近づく。ゆえに、任意に確率の低い状態は任意に小さな寄与をなし、厳密に不可能な状態は何も寄与しない。
\lecturetimestamp{01:18:27--01:19:46}}
\end{classroomqa}

\(M\)個の占有状態のそれぞれについて、\(p_i=1/M\)である。したがって
\begin{align}
 S
 &=-\sum_{\mathrm{occupied}\ i}
 \frac1M\log\!\left(\frac1M\right)\\
 &=-M\frac1M\log\!\left(\frac1M\right)\\
 &=-\log\!\left(\frac1M\right)
 =\log M.
 \label{eq:sm1-entropy-check}
\end{align}
一般式は特別な式と厳密に一致する。より広い分布はより多くの選択肢に無視できない重みを与え、したがってより大きなエントロピーを持つ。より狭い分布はより小さなエントロピーを持つ。

\begin{figure}[htbp]
\centering
\includegraphics[width=0.92\textwidth]{lecture_01_figure_06.png}
\caption{マークされた幅\(M\)は、有意な確率の広がりとエントロピーの間の関係を強調する。}
\lectureframe{01:16:42}
\end{figure}

エントロピーは状態空間上の確率分布に属する。それは一つの微視的状態のエネルギーや運動量のような裸の観測量ではない。厳密な状態はその状態に集中した分布で表され得て、そのときエントロピーはゼロであるが、知識の状態が変わればエントロピーは変化する。

\section{コイン、ビット、明示的な数え上げ}

\(n\)をコインの枚数とする。微視的配置の数は
\begin{equation}
 N=2^n.
 \label{eq:sm1-coin-state-count}
\end{equation}
である。コインについて何も知られていなければ、\(2^n\)個の配置はすべて等確率であり、
\begin{equation}
 S_{\max}=\log(2^n)=n\log2.
 \label{eq:sm1-coin-max-entropy}
\end{equation}
となる。エントロピーは独立なコインの枚数について加法的である。

\(2\)を底とする対数では、完全に未知の二状態系はエントロピー1ビットを持つ。自然対数では、エントロピー\(\log2\)を持つ。単位は慣習的なものであり、根底にある順序と加法性は変わらない。

もう一方の極端では、すべてのコインが既知ならば、ちょうど一つの配置が確率1を持つ：
\begin{equation}
 S=0.
\end{equation}
より多くの知識はより少ないエントロピーを意味する。

次に、\(n\)枚のコインのうちちょうど一枚が裏で、他はすべて表であり、裏の位置が未知であると仮定する。裏の可能な位置ごとに一つずつ、\(n\)個の等確率な配置があるので、
\begin{equation}
 S=\log n.
 \label{eq:sm1-one-tail-entropy}
\end{equation}
となる。これは一般に\(\log2\)の整数倍ではない。ビットは便利な単位であり、すべてのエントロピーが整数ビット刻みで生じるという主張ではない。

\begin{classroomqa}
\audiencequestion{\(\log2\)が現れるのは、各コインが二つの可能な状態を持つからだけなのか？三状態の要素なら\(\log3\)を与えるのか？}
\lecturerresponse{その通り。\(k\)個の可能な状態を持つ一様に未知の要素は\(\log k\)を寄与する。二値スイッチは、物理的スイッチが自然に二つの安定状態を実現するため、また二が最小の非自明な整数であるため、コンピュータにおいて特に便利である。
\lecturetimestamp{01:29:55--01:30:48}}
\end{classroomqa}

\begin{classroomqa}
\audiencequestion{対数にはどの底が意図されているのか？}
\lecturerresponse{これらの物理の公式では、\(\log\)は自然対数を意味する。情報理論家はしばしば底\(2\)を用いるので、エントロピーは直接ビットで測られる。底を変えることはすべてのエントロピーを一つの固定された換算係数で掛け、物理的結論を何も変えない。
\lecturetimestamp{01:30:49--01:32:06}}
\end{classroomqa}

\section{連続相空間におけるエントロピー}

古典相空間\(\Gamma\)に戻る。確率密度が体積\(V_\Gamma\)の領域上で一様で外部でゼロならば、\(M\)個の等確率状態を数えることの類似は
\begin{equation}
 S=\log\!\left(\frac{V_\Gamma}{V_0}\right).
 \label{eq:sm1-continuous-uniform}
\end{equation}
である。ここで\(V_\Gamma\)は座標と運動量の完全な高次元空間における体積であり、通常の空間体積ではない。固定された体積\(V_0\)は、対数を無次元にする参照相空間セルである。\(V_0=1\)と置くと、板書で用いられる略記\(S=\log V_\Gamma\)を与える。

一般の確率密度\(P(q,p)\)に対して、離散和は相空間積分になる。規格化は
\begin{equation}
 \int_\Gamma d\Gamma\,P(q,p)=1,
 \label{eq:sm1-continuous-normalization}
\end{equation}
であり、エントロピーは
\begin{equation}
 S=-\int_\Gamma d\Gamma\,
 P(q,p)\log\!\bigl[P(q,p)V_0\bigr].
 \label{eq:sm1-continuous-entropy}
\end{equation}
である。密度は相空間のセルあたりの確率である。概念的には、\eqref{eq:sm1-continuous-entropy}は依然として確率がどれだけ広く分布しているかを測っており、一様な塊に対しては、基準セルを単位として測った占有体積の対数に帰着する。その固定されたセルを変えると、すべての連続エントロピーが同じ加法定数だけずれ、エントロピー差は変わらない。

\begin{classroomqa}
\audiencequestion{なぜ座標記号が\(x\)から\(q\)に変わったのですか。}
\lecturerresponse{物理的な変化はない。古典力学では座標に\(x\)または\(q\)を、共役運動量に\(p\)を用いることがよくある。記法の切り替えは単なるミスであり、相空間の構成は同じである。
\lecturetimestamp{01:37:10--01:37:39}}
\end{classroomqa}

温度とボルツマン分布は意図的に後回しにされる。次の課題は熱平衡を定義し、次にその確率分布を特定することである。

\section{相関、加法性、情報に関する質問}

\subsection{相関とは、一方の測定が他方の分布を変えることである}

独立なコインについて完全に無知である場合、すべての配置に等しい確率を割り当てる。すると、1枚のコインを測定しても他のどのコインについても何もわからない。それらの分布は変わらない。

片側テールのアンサンブルは異なる。測定したコインが裏なら、他のすべてのコインは確実に表である。測定したコインが表なら、唯一の裏は残りの\(n-1\)枚のコインの中になければならないので、残りの各位置は条件付き確率
\begin{equation}
\begin{aligned}
 p(&\text{tail at a specified remaining position}\\
   &\mid\text{measured head})=\frac{1}{n-1}.
\end{aligned}
 \label{eq:sm1-conditional-tail}
\end{equation}
をもつ。これが相関である。すなわち、1つの変数を知ることが、別の変数に割り当てられる確率分布を変える。

\begin{classroomqa}
\audiencequestion{エントロピーの例では、コイン間の独立性が仮定されていましたか。}
\lecturerresponse{完全な無知がすべての\(2^n\)個の配置に等しい確率で表された場合に限る。その場合、1枚のコインを測定しても他のコインの分布は変わらない。片側テールの例は明示的に相関しており、測定は他の場所で何が可能かを更新する。述べられた情報が、どの確率分布を用いているかを決める。
\lecturetimestamp{01:38:43--01:41:58}}
\end{classroomqa}

非連結サイクルの例も同じ論理形式をもつ。保存された\(+1\)または\(-1\)のラベルを測定すると、それ以降のすべての状態が1つのセクターに制限される。保存情報で条件付けると、許される分布が変わる。

エントロピーは無相関な系に対して加法的である。もし
\begin{equation}
 p_{ij}=p_i^{(A)}p_j^{(B)},
\end{equation}
ならば
\begin{align}
 S_{AB}
 &=-\sum_{i,j}p_i^{(A)}p_j^{(B)}
 \log\!\left[p_i^{(A)}p_j^{(B)}\right]\\
 &=S_A+S_B.
 \label{eq:sm1-entropy-additivity}
\end{align}
である。相関はこの因数分解を妨げ、対の独立な不確かさを減らす。

演習として、二値変数の列を考えよう。ある変数が上向き状態にあると、各隣接変数が下向きである確率が四分の三になるものとする。一つのサイトを測定すると、近傍のサイトについて知られていることが直ちに変化する。このような分布は相関をもち、そのエントロピーは、それらの相関を考慮に入れずに孤立した一サイトのエントロピーを足し合わせるだけでは求められない。

\subsection{ボルツマン、シャノン、量子不確定性}

\begin{classroomqa}
\audiencequestion{古い熱力学的エントロピーは、シャノンの情報理論的エントロピーとどのように関係していますか。}
\lecturerresponse{同じ関数形が両方に現れる。
\(-\sum_i p_i\log p_i\)。シャノンは通常底\(2\)を用いたが、物理学者は通常自然対数を用い、エントロピーを熱力学的単位で表すときにはボルツマン定数を戻す。数学的な違いは固定された換算係数だけであるが、確率を与える物理的な問いは大きく異なりうる。
\lecturetimestamp{01:44:05--01:46:06}}
\end{classroomqa}

等確率のミクロ状態に対して、熱力学的公式はしばしば
\begin{equation}
 S_{\mathrm{thermo}}=k_{\mathrm B}\log W,
 \label{eq:sm1-boltzmann-formula}
\end{equation}
と書かれる。ここで\(W\)は到達可能なミクロ状態の数である。\(k_{\mathrm B}=1\)とする単位では、これはまさに\(S=\log W\)である。ボルツマン定数は温度の単位とエネルギーの単位を換算する。温度の自然な微視的スケールはエネルギーである。

\begin{classroomqa}
\audiencequestion{このエントロピーは単にハイゼンベルクの不確定性原理の別の形ですか。}
\lecturerresponse{いいえ。ここでのエントロピーは、可能な状態上の確率分布における不確かさを測る。量子論の言葉では、それは混合状態の不確かさであり、純粋状態でも残りうる本質的な測定の不確かさではない。この2つの概念は区別して保たなければならない。
\lecturetimestamp{01:46:07--01:47:31}}
\end{classroomqa}

\section{確立されたこと}

出発点は法則におけるランダム性ではなかった。決定論的力学は、初期状態に関する不完全な知識と結び合わされるだけで、すでに確率を生み出す。単一の軌道は各状態に費やされる時間を通じて重みを割り当てる。互いに連結していない軌道は保存量を明らかにし、各セクター間の事前重みを必要とする。許容される微視的法則は区別を保ち、有限状態空間ではこれが各状態における一本の入る矢印と一本の出る矢印を意味し、連続力学ではリウヴィルの定理となる。

その構造はまず、\(M\) 個の等確率で占有された状態に対して
\begin{equation}
 S=\log M
\end{equation}
を導き、次に
\begin{equation}
 S=-\sum_i p_i\log p_i,
\end{equation}
へ、そして最終的にその相空間積分へと導く。エントロピーは確率分布を、したがって知識の状態を記録する。細粒情報は、実用的な粗粒化記述がその追跡を見失う場合でも、微視的法則によって保存される。

エネルギー保存は第一法則を与えるが、部分系のエネルギーの加法性は相互作用エネルギーが無視できることを要求する。エントロピーは並列的な条件の下で加法的である。すなわち、部分系の確率が無相関でなければならない。これら二種類の加法性が次の段階を準備する。そこでは、熱平衡、温度、ボルツマン分布がエネルギーを確率へと結びつけるであろう。